
A four-hypothesis validation study tests the feasibility of zero-annotation semantic validation for research pipelines using MCP tool-calling traces. The experimental design follows an incremental validation strategy where each hypothesis verifies one step: trace completeness (H-E1), natural language presence (H-M1), semantic extraction quality (H-M2), and constraint inference (H-M3).

\textbf{Dataset.} Twenty real MCP trace logs from research pipeline executions (10 successful, 10 failed) provide balanced ground truth for failure detection validation. Each trace captures a complete pipeline run, logging all tool calls with their parameters and results in JSONL format. The dataset contains 596 tool calls across research tools (literature search, hypothesis generation, experiment design) and data processing tools (trace parsing, metric calculation, validation). Two known failure cases with documented reasoning errors serve as ground truth: h-e1 (data validation task) and h-m1 (effective rank contradiction where hypothesis assumed rank would decrease but experiments showed 6.02\% increase).

\textbf{Experimental Protocol.} H-E1 (Trace Completeness) validates that $\geq 95$\% of tool calls contain all required fields with natural language content ($\geq 10$ words). This threshold establishes the foundation---without complete traces, semantic analysis cannot extract assumptions or claims.

H-M1 (Natural Language Presence) measures the percentage of tool calls with $\geq 10$ words of text in query parameters or result content, targeting $\geq 90$\% presence. This validates the assumption that MCP traces encode reasoning in natural language.

H-M2 (Semantic Extraction Quality) uses a 50-sample validation set (25 queries, 25 results) with independent human annotation as ground truth. Zero-shot LLM extraction (temperature 0.0) with multi-vote consensus (3 votes, $\geq 2/3$ threshold) is evaluated against human annotations. Success requires $\geq 80$\% recall (finds most key items), $\geq 70$\% precision (low hallucination rate), and $\geq 70$\% inter-rater agreement (Cohen's kappa, validates gold standard reliability).

H-M3 (Constraint Inference) compares 8 assumptions extracted from early-phase queries against 150 claims from later-phase results, generating 1,200 assumption-claim pairs. Sentence-transformer embeddings with cosine similarity threshold <0.3 flag contradictions. Ground truth comes from the known h-m1 failure. The target is $\geq 70$\% recall with acceptable threshold at $\geq 60$\%.

\textbf{Evaluation Metrics.} Recall = TP / (TP + FN) measures percentage of actual failures detected. Precision = TP / (TP + FP) measures percentage of flagged violations that are real failures. Cohen's Kappa ($\kappa$) measures inter-rater agreement between LLM extraction and human annotation, controlling for chance. $\kappa \geq$ 0.70 indicates substantial agreement. False Positive Rate = FP / (FP + TN) is limited to <30\% to avoid excessive false alarms.

\textbf{Baselines.} Random prediction baseline randomly flags tool call pairs as contradictory with probability matching the dataset's failure rate (50\%), establishing the floor with expected precision/recall $\approx$ 50\%. Syntactic-only baseline validates traces using only JSON Schema checks (required fields, type constraints), representing existing tools like MLflow and DVC that focus on structural correctness without semantic analysis.
